\chapter{Radicals and classical applications}\label{chap:M07.15}
The polynomial Galois group $\Gal(f/K)$ means the automorphism group over $K$ of its splitting field. This notation is meaningful even for an inseparable polynomial; the extension itself is called Galois only when it is also separable. All root fields are embedded in a specified algebraic closure. In characteristic zero, a complex realization uses the fundamental theorem of algebra as an explicitly prior complex-analysis theorem.

% Topic grouping: Solvable groups and roots of unity
\KEntry{K07.15.solvable_groups}{Solvable groups}
\begin{definition}
Set $G^{(0)}=G$ and $G^{(j+1)}=[G^{(j)},G^{(j)}]$. The group is solvable if $G^{(r)}=1$ for some finite $r$.
\end{definition}
\begin{theorem}
For finite groups this is equivalent to a subnormal series with cyclic prime-order factors. Abelian groups are solvable; $S_5$ and every $S_n$ with $n\geq5$ are not.
\end{theorem}
\begin{proof}
Each derived subgroup is characteristic in its predecessor and has abelian quotient, since a quotient is abelian exactly when it kills every commutator. In a finite abelian group choose a subgroup of prime order using Cauchy's theorem, and induct on its quotient order. Pulling back its quotient series gives a series with cyclic prime-order factors. Refine each abelian quotient of the terminating derived series this way. Conversely a subnormal series with abelian factors gives solvability by the extension property proved immediately below, inductively from the trivial group. This argument uses that property only for the converse; its proof does not use the series characterization.

An abelian group's first derived subgroup is trivial. The earlier alternating-group theorem (Entry~\ref{K07.06.simple_groups_boundary}) proves $A_5$ is nonabelian simple. Its commutator subgroup is a nontrivial normal subgroup, so equals $A_5$; its derived series never terminates. Since $A_5\leq S_5\leq S_n$, the subgroup property below proves the nonsolvability assertions. Subnormal means normal in the next member, not necessarily normal in all of $G$.
\end{proof}

\KEntry{K07.15.solvable_closure}{Subgroups, quotients, and extensions}
\begin{theorem}
Subgroups and quotients of solvable groups are solvable. If $N\trianglelefteq G$ and both $N,G/N$ are solvable, then $G$ is solvable.
\end{theorem}
\begin{proof}
For a subgroup $H$, commutators computed in $H$ are the same elements as in $G$, so $H^{(j)}\subseteq G^{(j)}$ inductively. A surjection $q:G\to Q$ maps the subgroup generated by commutators onto the subgroup generated by their images; hence $q(G^{(j)})=Q^{(j)}$ by induction. This proves quotient closure. If $(G/N)^{(r)}=1$, then $G^{(r)}\subseteq N$. If $N^{(s)}=1$, repeated commutators give $G^{(r+s)}\subseteq N^{(s)}=1$, proving the extension property. Consequently any finite subnormal series with solvable factors makes its total group solvable.
\end{proof}

\KEntry{K07.15.roots_unity}{Roots of unity}
\begin{theorem}
For $n\geq1$ with $n\ne0$ in $K$, the roots $\mu_n$ of $X^n-1$ in an algebraic closure form a cyclic group of order $n$. If $\zeta$ generates it, the primitive roots are $\zeta^a$ with $\gcd(a,n)=1$. The extension $K(\mu_n)/K$ is finite Galois with abelian group embedded in $(\Z/n\Z)^\times$.
\end{theorem}
\begin{proof}
The derivative $nX^{n-1}$ has no common root with $X^n-1$, so the polynomial has $n$ distinct roots. Products and inverses of roots are roots; the finite multiplicative-subgroup lemma makes this group cyclic. The order-of-a-power formula proves the primitive-root criterion. All roots are powers of $\zeta$, so $K(\mu_n)=K(\zeta)$ is the separable splitting field of $X^n-1$. Any automorphism sends $\zeta$ to $\zeta^a$ of the same order, with $a$ a unit modulo $n$, and is determined by that value. Composition multiplies exponents because a field homomorphism preserves integer powers. This is an injective group homomorphism into the abelian unit group. If the characteristic is $p\mid n$, the distinct-root assertion fails: $X^p-1=(X-1)^p$.
\end{proof}

\KEntry{K07.15.cyclotomic_polynomials}{Cyclotomic polynomials}
\begin{theorem}
For $n\geq1$ let $\Phi_n(X)$ be the product of $X-\zeta$ over primitive complex $n$th roots. Then $\Phi_n\in\Z[X]$ is monic, has degree $\varphi(n)$, satisfies $X^n-1=\prod_{d\mid n}\Phi_d(X)$, and is irreducible over $\Q$. It is the minimal polynomial of every primitive $n$th root over $\Q$.
\end{theorem}
\begin{proof}
The roots of $X^n-1$ are partitioned according to their exact orders $d\mid n$; this proves the product identity, and counting primitive elements of the cyclic root group gives degree $\varphi(n)$. Inductively suppose all $\Phi_d$ for proper divisors are in $\Z[X]$. Divide $X^n-1$ by their monic product within $\Z[X]$. The complex factorization shows the remainder is zero and the quotient is $\Phi_n$. This proves integrality, starting with $\Phi_1=X-1$.

Let $f$ be the monic rational minimal polynomial of a primitive root $\zeta$. Gauss's lemma makes $f\in\Z[X]$ and a factor of $\Phi_n$. We claim that whenever $\ell$ is prime with $\ell\nmid n$, $\zeta^\ell$ is also a root of $f$. Otherwise let $g\ne f$ be its monic irreducible factor in $\Phi_n$. Then $f\mid g(X^\ell)$ in $\Z[X]$, by evaluation and monic division. Reduction modulo $\ell$ gives $\overline f\mid \overline g(X)^\ell$. A monic irreducible factor $h$ of the nonconstant $\overline f$ then divides $\overline g$ as well. But $fg\mid\Phi_n$, so $h^2$ divides $\overline\Phi_n$ and hence $X^n-1$ over $\F_\ell$. This contradicts its separability, since $\ell\nmid n$. The claim applies to any primitive root of $f$, so iteration over the prime factors of every positive $a$ coprime to $n$ makes every $\zeta^a$ a root of $f$. These are all $\varphi(n)$ primitive roots. Therefore $f=\Phi_n$. This proves irreducibility over $\Q$, and Gauss's lemma gives its integral primitive-polynomial formulation as well.
\end{proof}

\begin{remark}
Rational cyclotomic irreducibility is not irreducibility over every coefficient field. The factorization $X^4-1=(X-1)(X+1)(X^2+1)$ gives $\Phi_4=X^2+1$, which is irreducible over $\Q$ by the theorem. In $\F_5[X]$, however,
\[\Phi_4(X)=X^2+1=(X-2)(X+2),\]
since $-4=1$ in $\F_5$. Both displayed factors have degree $1$, so this is a genuine reducibility example even though the characteristic $5$ does not divide $4$.
\end{remark}

\KEntry{K07.15.cyclotomic_galois}{Rational cyclotomic Galois groups}
\begin{theorem}
For a primitive $n$th root $\zeta_n$,
\[[\Q(\zeta_n):\Q]=\varphi(n),\qquad
\Gal(\Q(\zeta_n)/\Q)\simeq(\Z/n\Z)^\times,
\]
where the unit class $a$ corresponds to $\zeta_n\mapsto\zeta_n^a$.
\end{theorem}
\begin{proof}
Cyclotomic irreducibility gives the degree. This field is the separable splitting field of $X^n-1$. The exponent homomorphism into the unit group was proved injective above. Both groups have cardinality $\varphi(n)$, so it is surjective. Thus every displayed exponent substitution really extends to an automorphism. For $n=1,2$ both sides are trivial. Over an arbitrary base field only the injection is automatic, not surjectivity onto all units.
\end{proof}

% Topic grouping: Radical towers and the two implications
\KEntry{K07.15.radical_definition}{Radical towers}
\begin{definition}
A radical tower over $K$ in an algebraic overfield is a finite tower $K=K_0\subseteq\cdots\subseteq K_r$ with $K_{i+1}=K_i(a_i)$ and $a_i^{n_i}\in K_i$ for positive integers $n_i$. An element is solvable by radicals if it belongs to some such tower. Repeated or degree-$1$ steps are permitted.
\end{definition}
\begin{lemma}\label{lem:radical-composita}
The compositum of finitely many radical towers in one overfield is radical. A base-preserving embedding maps a radical tower to a radical tower. The normal closure of a finite radical extension is a finite normal radical extension, without any separability assumption.
\end{lemma}
\begin{proof}
For two towers first adjoin all generators of the first, then those of the second. At a second-tower step, $b_i^{m_i}$ belongs to its preceding field, hence to the larger field now present; so the concatenated list remains radical. Induct for finitely many towers. Embeddings preserve every displayed power relation and field generation.

A radical tower is finite by the simple-degree and tower theorems. It has finitely many embeddings into the algebraic closure, by the embedding upper bound. Form the compositum $N$ of their image fields. This is finite and radical by the preceding observations. For a finite generating list of the original field, every root of every minimal polynomial occurs in one image: the simple embedding sending the generator to that root extends to the whole field. Thus $N$ is exactly the splitting field of the product of those minimal polynomials. It is normal, whether or not they are separable. This also proves the asserted normal-closure property.
\end{proof}

\KEntry{K07.15.polynomial_radicals}{Polynomial solvability by radicals}
\begin{definition}
A nonzero $f\in K[X]$ is solvable by radicals when all its roots in a fixed algebraic closure are radical over $K$.
\end{definition}
\begin{proposition}
This is equivalent to all roots lying in a common finite radical tower, and to its splitting field admitting a $K$-embedding into a finite radical extension.
\end{proposition}
\begin{proof}
There are finitely many distinct roots. If each has a radical tower, concatenate these towers using Lemma~\ref{lem:radical-composita}; the resulting field contains all roots and their generated splitting field. Conversely a common tower makes each root radical. If an abstract splitting field embeds into a radical extension, its image still contains every root of $f$ there: mapping a complete factorization supplies all factors. To compare with the fixed closure, extend the base embedding of that finite radical extension into the closure; its image is again a radical tower. This proves both representation-independent formulations and does not assume that having one radical root suffices for reducible $f$.
\end{proof}

\KEntry{K07.15.kummer_forward}{A radical step with roots of unity}
\begin{theorem}
Suppose $n\geq1$, $n\ne0$ in $K$, $\mu_n\subseteq K$, $a\in K^\times$, and $u^n=a$. Then $K(u)/K$ is cyclic Galois and its degree divides $n$.
\end{theorem}
\begin{proof}
All roots of $X^n-a$ are $\zeta u$, $\zeta\in\mu_n$, and lie in $K(u)$. They are distinct by the derivative criterion, since $a\ne0$ and $n\ne0$ in $K$. Thus $K(u)$ is its separable splitting field. The map
\[\Gal(K(u)/K)\longrightarrow\mu_n,\qquad\sigma\longmapsto\sigma(u)/u
\]
is injective, since $u$ generates the extension. It is a homomorphism: for $\tau(u)=\eta u$, the factor $\eta\in K$ is fixed by $\sigma$, so $(\sigma\tau)(u)/u=\eta\sigma(u)/u$. Its image is a subgroup of the cyclic order-$n$ group, hence cyclic of order dividing $n$. The Galois degree formula finishes the proof. The binomial itself need not be irreducible. If $a=0$ the field is simply $K$, handled separately. Without roots of unity, $\Q(\sqrt[3]{2})/\Q$ is not normal.
\end{proof}

\KEntry{K07.15.kummer_converse}{A cyclic extension is a radical step}
\begin{theorem}
Let $L/K$ be cyclic Galois of finite degree $n$, with $n\ne0$ in $K$ and $\mu_n\subseteq K$. There exists $u\ne0$ with $L=K(u)$ and $u^n\in K$.
\end{theorem}
\begin{proof}
Choose a generator $\sigma$ of the group and a primitive $n$th root $\zeta\in K$. The $K$-linear operator
\[R(x)=\sum_{j=0}^{n-1}\zeta^{-j}\sigma^j(x)
\]
is not identically zero, by the linear independence of the distinct embeddings $\sigma^j$ (Lemma~\ref{lem:independent-embeddings}). Choose $x$ with $u=R(x)\ne0$. Reindexing the sum using $\sigma^n=1$ gives $\sigma(u)=\zeta u$. Consequently the $n$ values $\sigma^j(u)=\zeta^j u$ are distinct, so $u$ has degree at least $n$ and hence generates $L$. Its $n$th power is fixed by $\sigma$, and therefore by the whole Galois group, so belongs to $K$. For $n=1$ choose $u=1$. The root-of-unity hypothesis is essential; characteristic-$p$ cyclic degree-$p$ extensions are not supplied by this Kummer argument.
\end{proof}

\KEntry{K07.15.radicals_imply_solvable}{Radicals force a solvable Galois group}
\begin{theorem}
Over any field $K$, if $a$ is radical over $K$, the polynomial Galois group of $m_{a,K}$ is solvable. If an irreducible $f$ has one radical root, $\Gal(f/K)$ is solvable. More generally a polynomial all of whose roots are radical has solvable polynomial Galois group, even in positive characteristic.
\end{theorem}
\begin{proof}
Place the given radical element, or the finitely many given radical roots, in a finite radical tower $T/K$ and take its finite normal radical closure $N/K$ by Lemma~\ref{lem:radical-composita}. Write a radical generating tower for $N$ with exponents $n_i$. In characteristic $p>0$, write $n_i=p^{e_i}m_i$ with $p\nmid m_i$; in characteristic zero put $m_i=n_i$. Let $m$ be the least common multiple of the finitely many $m_i$ and adjoin a primitive $m$th root $\zeta$. Put $C=K(\zeta)$ and $N'=CN$. This is still finite and normal over $K$: each factor is a splitting field and their compositum splits the product of their defining polynomials. The extension $C/K$ has abelian automorphism group by the roots-of-unity theorem.

Base-change the radical tower to $C$: each step is $M(u)/M$ with $u^{n_i}=b\in M$. If $u=0$ the step is trivial. Otherwise all its distinct conjugate roots are $\eta u$ with $\eta\in\mu_{m_i}\subseteq C$. In characteristic $p$, this follows from $X^{p^{e_i}m_i}-b=(X^{m_i}-u^{m_i})^{p^{e_i}}$ in the overfield. Thus $M(u)/M$ is normal, though possibly inseparable. Its automorphism group injects homomorphically into the cyclic $\mu_{m_i}$ by the same ratio calculation as before; this calculation needs no separability.

Let $G=\Aut_K(N')$. First its subgroup fixing $C$ is normal with abelian quotient, by restriction to $C$. Then take the subgroups fixing the successive fields of the base-changed tower. At each stage restriction is defined because that step is normal; its kernel fixes the next field and its image is a subgroup of the cyclic step automorphism group. Hence the resulting finite subnormal series has abelian factors and ends at $1$. The extension property makes $G$ solvable.

The splitting field $S$ of the polynomial under consideration is contained in $N'$ and normal over $K$. Restriction $G\to\Aut_K(S)$ is onto: an automorphism of $S$ extends to an embedding of $N'$ into the closure, and normality of $N'/K$ makes its image $N'$. The polynomial group is therefore a quotient of $G$ and is solvable. For a single radical element use the splitting field of its minimal polynomial; all its conjugates lie in $N$. An irreducible polynomial with that root is a scalar associate of this minimal polynomial. No characteristic-zero assumption was introduced in this proof.
\end{proof}

\KEntry{K07.15.radical_equivalence_char_zero}{The characteristic-zero equivalence}
\begin{theorem}
For $\operatorname{char}K=0$ and $0\ne f\in K[X]$, $f$ is solvable by radicals if and only if $\Gal(f/K)$ is solvable.
\end{theorem}
\begin{proof}
The forward implication is the preceding theorem, including repeated factors. For the converse let $S/K$ be the splitting field. It is separable in characteristic zero and hence finite Galois. Choose a subnormal series $G=G_0\trianglerighteq\cdots\trianglerighteq G_r=1$ with cyclic prime-order factors, using solvability. Put $K_i=S^{G_i}$. Applying the correspondence to $S/K_i$ gives cyclic Galois $K_{i+1}/K_i$ of degree $d_i=|G_i/G_{i+1}|$. No assumption that every $G_i$ is normal in all of $G$ is needed.

Adjoin to $K$ all $d_i$th roots of unity, equivalently one primitive root of order their least common multiple, and call the resulting field $C$. This is itself a radical step, since that root has a positive power equal to $1$. In the tower $CK_i$, base change of a cyclic Galois step remains cyclic Galois, with degree $e_i$ dividing $d_i$, by the Galois base-change theorem. Its base contains $\mu_{e_i}$, and characteristic zero makes $e_i$ invertible. The cyclic-to-radical theorem supplies a generator $u_i$ with $u_i^{e_i}\in CK_i$. Thus $C=CK_0\subseteq\cdots\subseteq CK_r=CS$ is a radical tower. Prepending $K\subseteq C$ gives a common radical extension containing every root of $f$. Constants have no roots and splitting field $K$, so are included. The proof does not assert this converse over unrestricted positive-characteristic fields.
\end{proof}

\KEntry{K07.15.small_degrees}{Equations of degree at most four}
\begin{theorem}
Every nonzero polynomial of degree at most $4$ over a characteristic-zero field is solvable by radicals.
\end{theorem}
\begin{proof}
The Klein group $V=\{1,(12)(34),(13)(24),(14)(23)\}$ is abelian and normal in $S_4$, since conjugation permutes these three double transpositions. The series $1\trianglelefteq V\trianglelefteq A_4\trianglelefteq S_4$ has abelian factors of orders $4,3,2$, so $S_4$ is solvable. For a polynomial with repeated factors, first take the product of its distinct monic irreducible factors: in characteristic zero this is separable, has degree at most four, and has the same splitting field. A separable polynomial with at most four distinct roots has Galois group embedded in $S_4$ (extend its root permutations by fixed points if necessary). Subgroups are solvable, so the radical equivalence applies.

For explicit formulas, completing the square gives the quadratic formula when $2$ is invertible. Replacing $x$ by $y-a/3$ in $x^3+ax^2+bx+c$ gives $y^3+py+q$, with $p=b-a^2/3$ and $q=2a^3/27-ab/3+c$. Choose $u^3=-q/2+\sqrt{q^2/4+p^3/27}$. If $u\ne0$, set $v=-p/(3u)$; multiplication of the two cubic radicands gives $v^3=-q/2-\sqrt{q^2/4+p^3/27}$. Then $u^3+v^3=-q$ and $3uv=-p$, so $(u+v)^3+p(u+v)+q=0$. The three roots are $\zeta^j u+\zeta^{-j}v$, $0\leq j<3$, counted with multiplicity. If the chosen $u$ is zero, the radicand product forces $p=0$, and cube roots of $-q$ supply all roots directly.

For a depressed quartic $y^4+py^2+qy+r$, if $q=0$ solve the quadratic in $y^2$. Otherwise choose a root $z$ of the cubic $q^2=4(2z-p)(z^2-r)$. Then $2z-p\ne0$; choose $u^2=2z-p$ and put $v=-q/(2u)$. The cubic identity gives $v^2=z^2-r$, so
\[y^4+py^2+qy+r=(y^2+z)^2-(uy+v)^2.
\]
Two quadratics now supply all four roots. Removing a cubic term of a general quartic uses $x=y-a/4$. These verifications explain the denominators $2,3$ in the usual formulas and permit complex intermediate radicals even when a desired root is real.
\end{proof}
\KEntry{K07.15.generic_unsolvability}{The general polynomial}
\begin{theorem}
Let $K$ have characteristic zero and let $s_1,\ldots,s_n$ be algebraically independent over $K$. The polynomial
\[T^n-s_1T^{n-1}+s_2T^{n-2}-\cdots+(-1)^n s_n
\]
over $K(s_1,\ldots,s_n)$ has Galois group $S_n$. For $n\geq5$ it is not solvable by radicals.
\end{theorem}
\begin{proof}
Begin with independent indeterminates $x_1,\ldots,x_n$ and $L=K(x_1,\ldots,x_n)$. The symmetric group acts faithfully by permuting them: different permutations give different values at some $x_i$. Let $F=L^{S_n}$ and let $e_i$ be the elementary symmetric polynomials in the $x_j$. Put $E=K(e_1,\ldots,e_n)\subseteq F$. The polynomial $\prod_j(T-x_j)$ has coefficients in $E$, distinct roots, and splitting field $L$. Thus $[L:E]\leq n!$ by the factorial bound. Artin's theorem gives $[L:F]=n!$, so the tower law forces $F=E$ and $[L:E]=n!$. Its automorphism group is therefore exactly the given $S_n$.

We verify that $e_1,\ldots,e_n$ are independent, so the result really applies to an unrestricted generic coefficient tuple. The elementary exchange argument needed here is as follows. If $x$ is transcendental over a field $B$ but algebraic over $B(y)$, clearing denominators in an annihilating equation gives a nonzero $P(Y,X)\in B[Y,X]$ with $P(y,x)=0$. It must depend on $Y$, since $x$ is transcendental over $B$. The polynomial $P(Y,x)$ is nonzero, again by transcendence of $x$, and proves $y$ algebraic over $B(x)$.

Consequently if a field $M$ is algebraic over $K(b_1,\ldots,b_m)$, it cannot contain $m+1$ independent elements over $K$. To prove this, suppose independent $a_1,\ldots,a_{m+1}$ exist. At a stage with $M$ algebraic over $K(a_1,\ldots,a_j,\text{remaining }b)$, choose a minimal subset of the remaining $b$ making $a_{j+1}$ algebraic over its generated field and $K(a_1,\ldots,a_j)$. It is nonempty by independence. Remove one member $y$; minimality makes $a_{j+1}$ transcendental over the field generated by all other chosen members. The exchange argument replaces $y$ by $a_{j+1}$ while leaving $M$ algebraic over the new full generating list. After $m$ replacements, $M$ is algebraic over $K(a_1,\ldots,a_m)$, contradicting independence of $a_{m+1}$.

If the $e_i$ were dependent, a maximal independent subset of them would have size $m<n$. Each omitted $e_i$ is algebraic over that subset: a dependence relation becomes a nonzero polynomial in $e_i$ because the subset is independent. Thus $E$, and then its finite extension $L$, would be algebraic over a field with $m$ generators. This contradicts the just proved bound applied to the $n$ independent $x_j$. Hence the $e_i$ are independent. The isomorphism sending the prescribed $s_i$ to $e_i$ transports the polynomial and its splitting field, proving the asserted generic group. For $n\geq5$ this group is not solvable, so the radical implication rules out a radical solution. This conclusion concerns independent coefficients; it does not assert that every specialized degree-$n$ polynomial is insoluble.
\end{proof}

\KEntry{K07.15.explicit_S5_quintic}{A specific quintic with group \texorpdfstring{$S_5$}{S5}}
\begin{theorem}
The polynomial $f=X^5-6X+3\in\Q[X]$ has Galois group $S_5$, and none of its roots is expressible by radicals over $\Q$.
\end{theorem}
\begin{proof}
Eisenstein at $3$ gives irreducibility. Let $L\subseteq\C$ be its splitting field, using the prior fundamental theorem of algebra. The root degree is $5$, so $5\mid[L:\Q]=|G|$ and Cauchy's theorem supplies an element of order $5$, necessarily a $5$-cycle in the faithful root action. The values $f(-2)=-17$, $f(0)=3$, $f(1)=-2$, $f(2)=23$ give three different real roots in the three disjoint intervening intervals by the intermediate value theorem. There cannot be four real roots: Rolle's theorem would give at least three distinct real zeros of $f'=5X^4-6$, which has just two. Hence exactly three roots are real, and the other two form a conjugate pair. Complex conjugation preserves $L$ and acts on the five roots as their single transposition.

A $5$-cycle and any transposition generate $S_5$. To check this assertion, number the vertices modulo $5$ so the cycle is translation by $1$. If the transposition exchanges $i,j$, its cycle conjugates give all edges joining $k$ to $k+(j-i)$. This graph is connected since every nonzero difference generates the additive group modulo $5$. Edge transpositions of any connected graph generate all transpositions: along a path, repeatedly use $(a\ c)=(b\ c)(a\ b)(b\ c)$ to extend the available endpoint transposition. Thus they generate the full symmetric group. Both permutations are in $G$, so $G=S_5$. Its nonsolvability and the irreducible one-root radical implication show no root is radical. The argument specifies both the polynomial and the allowed operations, and says nothing against numerical approximation.
\end{proof}

% Topic grouping: Constructibility
\KEntry{K07.15.constructible_definition}{Initial data and allowed constructions}
\begin{definition}
Fix a finite set $P\subseteq\R^2$ containing $O=(0,0)$ and $I=(1,0)$. Its coordinate field is $K=\Q(x,y:(x,y)\in P)\subseteq\R$. A point is constructible from $P$ if it occurs after finitely many intersections of lines through distinct available points and circles with available center and radius equal to a distance between available points. At each intersection either of the finitely many intersection points may be selected. Coincident curves are excluded. The straightedge is unmarked.
\end{definition}
The coordinate field includes both coordinates, not merely the field generated by a complex encoding of $P$ without its conjugates. Initial directions are fixed by the two coordinate axes, the second obtained as a perpendicular through $O$.

\begin{lemma}\label{lem:geometric-field-operations}
From constructible real coordinates one can construct their negatives, sums, products, quotients with nonzero denominator, and nonnegative square roots. Coordinates of an available point can be projected onto the coordinate axes.
\end{lemma}
\begin{proof}
We first justify the elementary auxiliary directions. Equal-radius circles centered at distinct $A,B$, with radius $|AB|$, meet in two points; the line through them is the perpendicular bisector of $AB$, since subtracting their two equal-distance equations gives the perpendicular-bisector equation. To construct a perpendicular to a line through a point on it, intersect the line with a circle centered at that point, obtaining two equidistant endpoints, and use their perpendicular bisector. For a point off a given line, a circle centered at it with radius twice its distance to some available point on the line meets the line twice: that radius exceeds its perpendicular distance. The perpendicular bisector of the resulting chord passes through the given center. Doubling a positive length uses two consecutive compass copies on a line. Two perpendicular constructions give a parallel through a prescribed point. These operations use only allowed intersections. Orthogonal projections now construct the two coordinates of a point; compass transfer moves a signed coordinate between the axes, choosing the appropriate one of the two intersections.

For numbers on the horizontal axis, copying a signed length from the endpoint of another constructs sums and negatives. For products, construct $(a,0),(0,1),(0,b)$ and the parallel through $(0,b)$ to the line joining $(0,1)$ to $(a,0)$. Its horizontal intercept is $(ab,0)$: its parametric coordinates are $(ta,b-t)$ and the intercept has $t=b$. For $a/b$, $b\ne0$, use instead the parallel through $(0,1)$ to the line joining $(0,b)$ to $(a,0)$; its intercept has coordinates $(a/b,0)$. Vertical and zero cases satisfy the same line equations and yield zero when appropriate.

For $a>0$, draw the circle with diameter endpoints $(-1,0),(a,0)$. Its center is $((a-1)/2,0)$ and radius $(a+1)/2$, both constructible by the operations already obtained. At its intersections with the vertical axis the equation becomes $y^2=a$. Transfer the nonnegative coordinate to the horizontal axis. For $a=0$ use the origin. This proves each operation explicitly by the permitted geometry.
\end{proof}

\KEntry{K07.15.quadratic_tower_iff}{Quadratic towers and constructibility}
\begin{theorem}
A real number $a$ is constructible from $P$ if and only if it lies in the top of a finite tower of real fields $K=K_0\subseteq\cdots\subseteq K_r\subseteq\R$ with $[K_{i+1}:K_i]\leq2$. A point is constructible if and only if both coordinates lie in one such tower.
\end{theorem}
\begin{proof}
For necessity, a line through two points over a field $M\subseteq\R$ has an equation $Ax+By+C=0$ with coefficients in $M$. A circle with center and defining radius from such points has equation $x^2+y^2+Dx+Ey+F=0$ with coefficients in $M$, since its squared radius is a sum of two squared coordinate differences. Two nonparallel lines have their intersection in $M^2$ by solving the two linear equations. Substituting a line equation in a circle equation gives a polynomial of degree at most $2$ for one coordinate, with the other rationally determined; its real roots lie in $M(\sqrt d)$ for a nonnegative discriminant $d\in M$. Two distinct intersecting circles give a nonconstant line equation when their equations are subtracted, and reduce to this case. Concentric unequal circles have no intersection; coincident curves are excluded. Tangency can give a degree-$1$ step. Adjoining the needed square root at each of finitely many construction steps gives a tower as asserted, containing all produced coordinates.

For sufficiency, the lemma constructs every element of the initial coordinate field: all its generators are obtained by projection, and each rational expression uses finitely many permitted operations. In a nontrivial real quadratic extension choose $u$ outside the base field $M$. Its minimal polynomial is $X^2+bX+c$; completing the square gives $u=(-b\pm\sqrt{b^2-4c})/2$, with positive discriminant because the roots are real and distinct in characteristic zero. Hence the lemma constructs $u$ from available coefficients. Every element of $M(u)$ is $s+tu$ with $s,t\in M$ and is therefore constructible. Induct through the finite tower. Coordinates can be combined into a point by perpendicular lines. Conversely projecting a constructed point constructs each coordinate. If the two coordinates are obtained in different towers, concatenate their quadratic generators: after enlarging the preceding field each generator still satisfies its old quadratic equation, so its new step has degree at most $2$. This produces a common tower.
\end{proof}

\begin{lemma}\label{lem:two-group-chains}
If $G$ is a finite group of order a power of $2$ and $H\leq G$, there is a chain $H=H_0< H_1<\cdots<H_s=G$ with each index $[H_{i+1}:H_i]=2$.
\end{lemma}
\begin{proof}
Every proper subgroup $H$ has a strictly larger normalizer. Prove this by induction on $|G|$. The center $Z$ is nontrivial by the finite $p$-group center theorem, Entry~\ref{K07.06.pgroup_center}. If $Z\nsubseteq H$, the larger group $HZ$ normalizes $H$. If $Z\subseteq H$, apply induction in the smaller group $G/Z$ to $H/Z$; a larger normalizer there lifts to an element outside $H$ normalizing $H$. Thus $N_G(H)/H$ is a nontrivial $2$-group. Cauchy's theorem supplies a subgroup of order $2$ in this quotient; its inverse image is an overgroup of $H$ of index $2$ over $H$. Repeat with that overgroup. Each increase is strict and the finite process reaches $G$.
\end{proof}

\KEntry{K07.15.normal_closure_criterion}{The exact normal-closure criterion}
\begin{theorem}
For algebraic real $a$ over $K\subseteq\R$, $a$ is constructible over $K$ if and only if the normal closure $N/K$ of $K(a)/K$ has degree a power of $2$. Constructibility over $K$ means the quadratic-tower criterion; for finitely generated coordinate fields it agrees with the preceding geometric definition.
\end{theorem}
\begin{proof}
Suppose first that $a$ is in a finite real quadratic tower. Write each nontrivial step as $E_i=E_{i-1}(\sqrt{b_i})$; this follows by completing the square. Construct finite normal extensions $N_i/K$ of $2$-power degree containing $E_i$, starting with $N_0=K$. If $N_{i-1}$ is already constructed, let
\[N_i=N_{i-1}(\sqrt{\sigma(b_i)}:\sigma\in\Gal(N_{i-1}/K)).\]
Each of the finitely many adjunctions has degree at most $2$, so the new degree is a power of $2$. Any $K$-embedding preserves $N_{i-1}$ and permutes the displayed square roots up to signs, so preserves $N_i$. Hence $N_i/K$ is normal and, in characteristic zero, Galois. It contains $E_i$ by the identity conjugate. The normal closure of $K(a)$ is contained in the last $N_i$, so its degree divides a power of $2$ and is a power of $2$.

Conversely let $G=\Gal(N/K)$ be a $2$-group and $H=\Gal(N/K(a))$. Apply Lemma~\ref{lem:two-group-chains} and take fixed fields in reverse order. Consecutive degrees are $2$, and the fields run from $K$ to $K(a)$. All are contained in the real field $K(a)$, even if the normal closure itself has nonreal elements. Thus this is a real quadratic tower containing $a$, proving constructibility. This subgroup-relative chain is why a complex normal closure causes no real-coordinate obstruction.
\end{proof}

\KEntry{K07.15.degree_obstruction}{Necessary degree and its failed converse}
\begin{theorem}
If real algebraic $a$ is constructible over $K$, then $[K(a):K]$ is a power of $2$. The converse is false, even for degree $4$ over $\Q$.
\end{theorem}
\begin{proof}
A quadratic tower has total degree a power of $2$. The simple degree divides that total degree by the tower law, so is a power of $2$.

For a fully specified counterexample take a real root $a$ of $h=X^4-X-1$. Reduction modulo $2$ gives $X^4+X+1$, with no root in $\F_2$. The only monic irreducible quadratic over $\F_2$ is $X^2+X+1$ (checking the two possible coefficients proves this), and the remainder of $X^4+X+1$ modulo it is $1$. A reducible quartic without linear factors must have quadratic factors; thus the reduction is irreducible and so is $h$ over $\Q$. It has two real roots: signs at $-1,0$ and at $1,2$ give at least two, while three real roots would, by Rolle's theorem, give two zeros of $h'=4X^3-1$, which has only one. Therefore complex conjugation on the splitting field $L$ acts as a transposition on the four roots.

Number the roots $a_1,\ldots,a_4$ and form
\[b_1=a_1a_2+a_3a_4,\quad b_2=a_1a_3+a_2a_4,\quad b_3=a_1a_4+a_2a_3.
\]
Their differences are products of distinct root differences; for example $b_1-b_2=(a_1-a_4)(a_2-a_3)$, so they are distinct. Writing $e_i$ for the elementary symmetric functions of the roots, expansion gives
\[\sum b_i=e_2,\quad\sum_{i<j}b_i b_j=e_1e_3-4e_4,\quad
 b_1b_2b_3=e_1^2e_4+e_3^2-4e_2e_4.
\]
Here $e_1=e_2=0,e_3=1,e_4=-1$, so the $b_i$ are the roots of $Y^3+4Y-1$. Its only possible rational roots are $1,-1$, neither of which is a root; it is irreducible. Hence $3\mid[L:\Q]$. Also $4\mid[L:\Q]$ because $a_1$ has degree $4$. The faithful Galois group is a subgroup of $S_4$ with order divisible by $12$, so has order $12$ or $24$. An index-$2$ subgroup of $S_4$ must be $A_4$: its quotient map to $C_2$ has the same value on all conjugate transpositions, and that value is nonidentity since transpositions generate $S_4$; its kernel is exactly the even permutations. Our group contains a transposition, so cannot have order $12$. Thus $G=S_4$ and $[L:\Q]=24$.

The normal closure of $\Q(a)$ is this splitting field, so the exact criterion just proved makes $a$ nonconstructible although $[\Q(a):\Q]=4$. This example also shows a degree divisor need not occur: if a quadratic field lay between $\Q$ and $\Q(a)$, its fixing subgroup in $S_4$ would have order $12$ and contain the root stabilizer $S_3$. But the only order-$12$ subgroup is $A_4$, which cannot contain that $S_3$ with its transpositions. Thus the divisor $2$ of the root-field degree $4$ has no corresponding intermediate field.
\end{proof}
\KEntry{K07.15.double_cube}{Doubling a cube}
\begin{theorem}
From rational initial unit data no unmarked-straightedge-and-compass construction doubles the volume of the unit cube.
\end{theorem}
\begin{proof}
The required positive side length is $a=\sqrt[3]{2}$. Eisenstein at $2$ makes $X^3-2$ irreducible over $\Q$, so its degree is $3$. Constructible real lengths over the rational coordinate field have power-of-two degree, a contradiction. This is an exact obstruction for the stated instruments and initial data, not an obstruction to approximate numerical lengths or to constructions with a marked ruler.
\end{proof}

\KEntry{K07.15.trisect_angle}{No universal angle trisection}
\begin{theorem}
There is no straightedge-and-compass construction trisecting every prescribed angle. In particular the constructible angle $60$ degrees cannot be trisected with these instruments from ordinary unit data.
\end{theorem}
\begin{proof}
An equilateral triangle constructs $60$ degrees from a unit segment. If a trisection existed, concatenating it with that construction would make $20$ degrees constructible over $\Q$, hence make $b=2\cos(20\text{ degrees})$ constructible. Taking real parts in $(\cos\theta+i\sin\theta)^3$ gives $2\cos3\theta=(2\cos\theta)^3-3(2\cos\theta)$. Thus $b^3-3b-1=0$. The only possible rational roots of this monic cubic are $1,-1$, and substitution excludes both. A reducible cubic over a field has a linear factor, so it is irreducible and $[\Q(b):\Q]=3$. This contradicts the degree obstruction. Some angles do admit a trisection, for example a straight angle can be divided into three $60$-degree angles; the theorem denies a construction for every input angle.
\end{proof}

\KEntry{K07.15.square_circle}{Squaring the circle}
\begin{theorem}
Starting with rational coordinates and a unit segment, the unit-radius circle cannot be squared by an unmarked straightedge and compass.
\end{theorem}
\begin{proof}
A square with area equal to the circle would have positive side $s=\sqrt\pi$. Every constructible real number over $\Q$ is algebraic by the quadratic-tower theorem. If $s$ were algebraic, its square would be algebraic by closure of algebraic elements under multiplication, contradicting the transcendence of $\pi$. The analytic transcendence input is proved separately in Theorem~\ref{thm:pi-transcendental}; it is not claimed to follow just from elementary extension degrees. Historically the authoritative theorem is F.~Lindemann, \emph{Ueber die Zahl $\pi$}, \emph{Mathematische Annalen} \textbf{20} (1882), 213--225, DOI \texttt{10.1007/BF01446522}. The present paragraph proves the algebraic implication from that transcendence theorem and the elementary area formula. Fixing the ordinary rational initial data matters: if $\sqrt\pi$ were provided as initial data, the question would be different.
\end{proof}

\KEntry{K07.15.regular_polygons}{The Gauss--Wantzel criterion}
\begin{theorem}
For $n\geq3$, a regular $n$-gon is constructible by straightedge and compass from ordinary unit data if and only if $\varphi(n)$ is a power of $2$. Equivalently
\[n=2^a p_1\cdots p_s\]
with $a\geq0$ and distinct Fermat primes $p_j=1+2^{2^{b_j}}$.
\end{theorem}
\begin{proof}
First normalize to the unit circle with a vertex at $1$. This is equivalent to the ordinary construction problem: a constructed regular polygon has a constructible circumcenter from perpendicular bisectors, a positive constructible radius from its coordinate distance, and a constructible rotation and scaling sending one vertex to $1$. These operations use field operations and square roots on coordinates. Conversely any circle with a constructible center and nonzero constructible radius is obtained by such a similarity. The successive unit-circle vertices are powers of $\zeta=e^{2\pi i/n}$, so constructing $\zeta$ constructs the whole polygon by complex multiplication, implemented by real field operations.

If $\zeta=x+iy$ is constructible, its real coordinates lie in a common real quadratic tower $M/\Q$. Then $\zeta\in M(i)$, whose degree over $\Q$ is a power of $2$. Thus $[\Q(\zeta):\Q]=\varphi(n)$ divides a power of $2$. Conversely suppose $\varphi(n)$ is a power of $2$. The field $N=\Q(\zeta,i)$ is finite Galois: it is the compositum of the two separable splitting fields $\Q(\zeta)$ and $\Q(i)$. Its degree is $\varphi(n)$ times $1$ or $2$, hence a power of $2$. The real elements
\[x=(\zeta+\zeta^{-1})/2,\qquad y=(\zeta-\zeta^{-1})/(2i)
\]
belong to $N$. Their normal closures over $\Q$ lie in $N$, so have $2$-power degrees. The exact normal-closure criterion makes both constructible, hence makes $\zeta$ and the polygon constructible. This proves sufficiency using a quadratic tower, not by reversing the merely necessary element-degree condition.

For the arithmetic form, write $n=2^a\prod p_j^{e_j}$ with odd distinct primes. Counting units modulo each prime power and using the coprime Chinese remainder theorem gives
\[\varphi(n)=2^{\max(a-1,0)}\prod_j p_j^{e_j-1}(p_j-1).
\]
Indeed $p^e-p^{e-1}$ residues modulo $p^e$ are not divisible by $p$, and unit choices multiply under CRT. For this positive integer to be a power of $2$, every $e_j=1$ and each $p_j-1=2^{c_j}$. If $c_j=2^{b_j}u$ with odd $u>1$, then $p_j=(2^{2^{b_j}})^u+1$ is divisible by the smaller integer $2^{2^{b_j}}+1$, because $A^u+1$ is divisible by $A+1$ for odd $u$. Primality forces $u=1$, giving a Fermat prime. Conversely these distinct factors make every factor of the totient a power of $2$. This proves the equivalence and includes, for example, the constructibility of the $3$-, $5$-, and $17$-gons. It makes no assertion that all Fermat primes have already been found.
\end{proof}


